% Compile with XeLaTeX from the repository root.
\documentclass[twocolumn]{article}
\usepackage[a4paper,margin=20mm]{geometry}
\usepackage{gradbars,booktabs}
\title{Compact Results in a Two-Column Paper}
\author{gradbars example}
\date{}
\gradbarscolumn{G}{max=100,width=22mm,precision=1,label width=3em}
\begin{document}
\maketitle
\section{Evaluation}
The table reports illustrative accuracy scores. All rows use the same scale
from zero to 100; the bars are placed directly in the numeric column.
\begin{table}[ht]
\centering
\caption{Accuracy (\%). Shared range: 0--100.}
\begin{tabular}{@{}lG@{}}
\toprule
Model & \multicolumn{1}{c}{Accuracy}\\
\midrule
Baseline & 72.4\\
Compact & 81.6\\
Proposed & 92.7\\
Unavailable & NA\\
\bottomrule
\end{tabular}
\end{table}
\section{Interpretation}
A longer bar indicates a larger score. Missing values are displayed separately
from zero. Column headings bypass numeric collection with multicolumn.
\newpage
\section{Uncertainty}
Intervals should be described in the caption. These illustrative intervals
are supplied directly; the package does not estimate uncertainty.
\begin{center}
\begin{tabular}{@{}lc@{}}
\toprule
Model & Score and interval\\
\midrule
A & \gradbar[width=22mm,error=4,label width=3em]{72}\\
B & \gradbar[width=22mm,error=2,label width=3em]{89}\\
\bottomrule
\end{tabular}
\end{center}
\end{document}
